

\subsection{Adapting Vision--Language Models Under Domain Shift}

Large pre-trained vision and vision--language models exhibit strong transfer across downstream recognition tasks, but this transfer is not equivalent to robustness under distribution shift. Models such as CLIP, SigLIP, ALIGN, and DINOv2 are commonly evaluated or adapted on fine-grained recognition, medical imaging, satellite imagery, textures, corruptions, and ImageNet-derived shifts, where performance can degrade substantially relative to in-domain benchmarks \citep{radford2021learning, zhai2023sigmoid, jia2021scaling, oquab2024dinov2, hendrycks2019benchmarking}. This has motivated a large literature on adaptation.

Existing adaptation methods differ in how directly they modify the model. Prompt-learning methods such as CoOp, CoCoOp, MaPLe, PromptSRC, and TAP learn task-specific context tokens while keeping most or all backbone parameters fixed \citep{zhou2022learning, zhou2022conditional, khattak2023maple, khattak2023promptsrc, ding2025tap}. Parameter-efficient fine-tuning methods such as LoRA and related low-rank adaptation techniques modify internal weights through compact trainable updates \citep{hu2022lora}. Adapter-based and visual-prompt methods provide alternative mechanisms for shifting intermediate representations or input-level tokens without full fine-tuning \citep{gao2024clipadapter, jia2022visual}. In domain generalization and robustness settings, empirical risk minimization, distributionally robust optimization, source-free adaptation, and test-time prompt tuning are also used to improve target-domain performance \citep{sagawa2020distributionally, singha2023ad, shu2022tpt}.

However, most of this literature evaluates adaptation primarily through output behavior: accuracy, zero-shot transfer, robustness, or calibration. These metrics answer whether adaptation improves performance, but not what adaptation does to the internal feature basis. In particular, it remains unclear whether adaptation reuses pre-trained semantic features, repurposes them, suppresses them, or induces new domain-specific features. This distinction is central for interpretability: two adapted models with similar accuracy may rely on very different internal mechanisms.

\subsection{From Output Explanations to Internal Representation Analysis}

A large body of explainability work explains model predictions by attributing outputs to input regions or pixels, for example through saliency maps, gradient-based attribution, Grad-CAM, integrated gradients, and relevance propagation methods \citep{selvaraju2017grad, sundararajan2017axiomatic, bach2015pixel}. These methods are useful for localizing influential image regions, but they provide limited information about the feature-level organization of the model. This limitation is especially important under domain shift, where failures may arise from non-local cues such as texture, acquisition artifacts, background correlations, or entangled feature reuse.

A complementary direction studies intermediate representations directly. Earlier work analyzed internal visual features through feature inversion, network dissection, neuron-level interpretation, and representation similarity \citep{mahendran2014understanding, bau2017network, durrani2020analyzing, csiszarik2021similarity}. For vision--language models, the shared image--text embedding space enables text-grounded interpretation of visual representations. Gandelsman et al. decompose CLIP-ViT image representations into contributions from layers, attention heads, and image tokens, and interpret these contributions using text directions in the shared embedding space \citep{gandelsman2024interpreting}. This line of work shows that internal components of VLMs can be assigned meaningful semantic roles.

Nevertheless, architectural decompositions and attribution-style analyses do not directly answer whether the underlying feature basis is disentangled. A head, layer, token, or direction may contribute to a prediction while still combining multiple unrelated concepts. This is particularly problematic when studying adaptation: an adapted model may preserve the same high-level architectural pathways while changing the feature-level dictionary used inside those pathways. Our work therefore focuses on feature-level interpretability rather than only output attribution or architectural attribution.

\subsection{Sparse Autoencoders as Feature-Level Interpretability Tools}

Sparse autoencoders (SAEs) have emerged as a tool for decomposing neural activations into sparse, reusable features. Motivated by the problem of superposition, SAEs learn an overcomplete dictionary that reconstructs intermediate activations while encouraging sparse feature usage \citep{elhage2022superposition, bricken2023monosemanticity, cunningham2024sparse}. Ideally, individual SAE features are more monosemantic than raw neurons or dense activation directions, making them easier to inspect, compare, and intervene on.

Recent work has extended SAEs beyond language models to vision models and vision--language encoders. Rao et al. use CLIP SAE dictionaries as task-agnostic concept bottlenecks across downstream tasks \citep{rao2024discover}. Karvonen et al. train SAEs on CLIP and DINOv2 activations and demonstrate visual features that support causal interventions across recognition and segmentation settings \citep{karvonen2025saevision}. Pach et al. show that SAE interventions on CLIP's vision encoder can steer multimodal language model behavior without modifying the language model itself \citep{pach2025monosemantic}. Other work studies sparse concept embeddings, steerability, and hierarchical or multi-granularity SAE variants for visual representations \citep{bhalla2024splice, joseph2025steering, zaremba2025msae, bussmann2025matryoshka}.

Patch-level SAEs are especially relevant for vision transformers because they operate over spatial tokens rather than only global image embeddings. PatchSAE enables localized concept extraction and can reveal whether image patches activate features corresponding to objects, textures, styles, or other visual concepts \citep{lim2024patchsae}. This makes patch-level SAE analysis a natural tool for studying domain shift, where the relevant changes may be local, global, or both.

\subsection{Interpreting Adapted Models and the Dictionary-Transfer Gap}

The closest work to ours is PatchSAE, which applies a base-model SAE to MaPLe-adapted CLIP and finds that prompt-based adaptation largely remaps existing concepts rather than learning new ones \citep{lim2024patchsae}. This is an important result because it suggests that, at least for prompt tuning on the studied benchmarks, adaptation may operate through concept remapping rather than feature creation. However, this also exposes a critical assumption: the SAE dictionary trained on the base model must remain a valid analysis basis after adaptation.

Our work studies precisely when this assumption holds and when it breaks. Unlike prior SAE analyses that reuse a single base-model dictionary as a fixed interpretability lens, we compare base SAEs and domain-adapted SAEs after stronger adaptation methods and larger domain shifts. This distinction matters because a base SAE can appear interpretable while failing to faithfully reconstruct or causally explain the adapted model's internal computation. In such cases, the issue is not only that the adapted model has changed; it is that the interpretability instrument itself may have gone out of distribution.

We therefore position our work as a feature-level study of adaptation, not merely an adaptation benchmark or an SAE visualization study. Robustness benchmarks measure performance under shift; adaptation papers optimize that performance; text-grounded decompositions identify architectural contributions; prior SAE work extracts sparse visual concepts from fixed or lightly adapted models. What remains underexplored is whether the sparse feature dictionary learned from a base model continues to faithfully explain models adapted to substantially different domains. We address this gap by evaluating reconstruction fidelity, feature sparsity, monosemanticity, feature drift, and causal faithfulness across model architectures, adaptation families, and domain shifts.